% newton_cabling_setup.tex -- IEEEtran-ready fragment (no preamble).
% % newton_cabling_setup.tex -- IEEEtran-ready fragment (no preamble).
% % newton_cabling_setup.tex -- IEEEtran-ready fragment (no preamble).
% % newton_cabling_setup.tex -- IEEEtran-ready fragment (no preamble).
% \input{newton_cabling_setup} or paste before \end{document}.
% Needs only packages already in the preamble.

\providecolor{cmured}{RGB}{196, 18, 48}

%=====================================================================
\section{Simulation: Deformable Cable Insertion}
%=====================================================================

To study insertion with a deformable object, we simulate ethernet
cabling: the robot grips a flexible cable $\sim$5\,cm behind the
connector, so the RJ45 plug dangles from its strain-relief boot and
must be levelled and seated into a panel-mount jack with
sub-millimetre clearance. This section summarizes the design of the
simulation and the empirical lessons that shaped it.

\subsection{Choice of Physics Engine}

We evaluated MuJoCo, Genesis, Isaac Sim, and Newton. MuJoCo and
Genesis could not resolve the fine cable/connector interaction, which
left the two NVIDIA-stack options. We chose Newton over Isaac Sim
because its VBD solver treats one-dimensional deformables ---
self-collision, bend/twist transport, two-way coupling to rigid
bodies --- as a first-class, GPU-batched capability, and because it
is a headless Python library rather than a GUI-centric Omniverse
application, which kept iteration on contact behaviour fast. The cost
is maturity: PhysX articulations are far more battle-tested, and we
spent roughly fifteen debugging iterations discovering undocumented
VBD constraints (fixed joints freeze deformable chains unless
merged away; actuated joints must be softened; the rest pose must be
updated via forward kinematics before the solver is constructed;
stable joint-limit damping is $\sim 10^{-4}$, four orders of
magnitude below the conventional default).
Since Newton is also being integrated into Isaac Lab as a swappable
backend, we do not view the choice as locking us out of that
ecosystem.

\subsection{Cable and Contact Modelling}

We model the cable at two fidelities. For teacher training, the cable
is a Cosserat rod of 10\,mm capsule segments (radius 3.25\,mm,
spanning 6\,cm behind to 5\,cm ahead of the grasp), with the
connector attached as a mesh shape on the front segment ---
attaching it as a welded rigid body destabilizes VBD. A PPO teacher
trained in this environment reaches 94.4\,\% held-success at the
hardest curriculum stage. For large-scale data generation we
substitute a rigid approximation: a single free body carrying the
cord capsule, the plug mesh, and a small box at the grasp point that
gives the flat fingerpads a roll-lock a smooth capsule lacks.

Contact uses narrow-band signed-distance fields built directly on
the CAD meshes ($k_e = 10^{6}$, contact gap 0.1\,mm, socket
$\mu = 0.5$, cable $\mu = 2.0$), avoiding convex decomposition
entirely. Two findings dominated the design.
\textcolor{cmured}{First}, contact damping must be zero: the default
($k_d = 100$) acts through SDF shell contacts that fire at several
millimetres of actual separation, injecting up to
$\sim$400\,rad/s of spurious angular velocity into the light
(6.9\,g) connector. \textcolor{cmured}{Second}, contacts must never
act on articulation bodies: every grasp variant that allowed this
diverged. We therefore disable finger-mesh collision and realize the
grasp with two kinematic pad proxies placed at the measured pad
faces, closed to a $\sim$1.5\,mm effective squeeze on the 6.5\,mm
cable.

\subsection{Robot Model}

The arm (StandardBots RO1 with an AG-145 gripper) is driven
kinematically: end-effector deltas ($\pm$2\,mm, $\pm 0.75^{\circ}$
per 60\,Hz step) are converted to joint targets by
damped-least-squares IK and interpolated across the eight physics
substeps, so the shallow grasp is never violated within a step. For
hardware-faithful data we replace direct IK with the real controller
chain --- jerk-limited trajectory retiming, 500\,Hz streaming,
10\,ms transport delay, a 2\,ms plant --- which reaches 100\,\% held
success once
replanning is seeded with velocity and acceleration state (the naive
rest-to-rest replan advanced only 0.12\,\% of commanded motion).

\subsection{Digital-Twin Geometry}

Geometry comes from the manufacturer CAD models of the plug and jack
(McMaster-Carr 9953K216 and 1422N17). Our first conversion, a 0.10\,mm
voxel remesh, silently destroyed the task --- the real plug--cavity
clearance is $\sim$0.11\,mm, comparable to the voxel pitch, so the
gap quantized to zero. We instead tessellate the STEPs directly; the
narrow-band SDF tolerates the resulting non-watertight meshes and
preserves the true clearance. We also removed the latch catch ledge
from the physics model after finding that a rigid connector jams on
it (the real part inserts only because the plastic flexes). To keep
physics and rendered appearance consistent, a scan-matched variant
registers a photogrammetry scan of the real cable+plug into the CAD
frame (Kabsch alignment, sub-micron residual), so the collision mesh
matches the splat the cameras see. For real-world collection we
print a single-body fixture that holds the jack rigidly, with a
simulated twin registered to the socket to within 0.5\,mm.

\subsection{Determinism}

Kinematic quantities (joint streams, wrist poses) reproduce
bit-exactly across runs; contact dynamics do not, because VBD's
contact reduction uses GPU atomics and the resulting contact
ordering perturbs the friction impulse on the light connector ---
same-seed episodes near the success boundary can flip. We therefore
validate calibration transfer on kinematic channels (exact to
$2.5\times10^{-14}$ degrees) and evaluate dynamic outcomes
statistically.

%=====================================================================
\section{Photoreal Observations via Gaussian Splatting}
%=====================================================================

\subsection{Rendering Architecture and Scene}

Observations are produced by compositing the Newton state into a
Gaussian-splat reconstruction of the real scene. An external
renderer receives per-object world transforms and camera parameters
each frame and returns RGB; the interface is physics-agnostic, so
the same renderer serves our Isaac/PhysX pipeline. The scene comprises 19 splats: table, jack,
connector head and tail, and the 15 robot links.

Because a deformable cannot be represented by a rigid splat, the
cable itself is not rendered: the connector head and boot splats
ride the simulated plug-face pose, the cord behind them is simulated
but unrendered, and the plug splat is cropped to the rigid extent
the physics actually tracks. The plug splat is baked by
\emph{reusing the physics registration}, guaranteeing that pixels
and collision geometry cannot drift apart; the jack splat is
registered by point-pair correspondence, which also resolves the
arbitrary SfM scale of the reconstruction (0.8193; 0.91\,mm surface
residual). Robot splats keep only degree-0 spherical harmonics, as
the renderer misinterprets higher orders.

\subsection{Cameras}

All views use the intrinsics of the real Intel RealSense D415 at
$640\times480$ ($f_x{=}605.78$, $f_y{=}605.29$, $c_x{=}326.14$,
$c_y{=}229.88$), so simulated and real observations share a
projection model; both are square-resized to $512^2$ identically. Up
to four views are rendered --- front, side, mirror, and a dynamic
wrist camera whose pose and FOV are taken from the calibrated camera
mount in the robot model, so the eye-in-hand view is rigid by
construction. We
verify simulator--renderer alignment by rendering the Newton scene
from identical cameras and diffing against the GS frames; this
requires reproducing the wrist camera's $\sim$8$^{\circ}$ roll,
which the stock debug camera cannot represent.

\subsection{Shadows, Compositing, and Randomization}

Gaussians emit baked radiance, so dynamic objects cast no shadows
and appear pasted onto the scene. We recover shadows with a two-pass
path-traced ratio mask over the Newton geometry: a pass with and a
pass without dynamic occluders, whose luminance ratio (clamped at 1)
multiplies the GS frame. All terms except occlusion cancel between
the passes, so no segmentation is needed and the reconstruction's
own baked shadows are not double-darkened. Frames are then
gamma-corrected (1.8) and cleaned by Difix3D before dataset
conversion; at evaluation, live camera frames pass through the same
Difix so the policy sees its training distribution.

Domain randomization spans both halves of the pipeline: background
yaw over $U(0,360^{\circ})$ and shadow strength over $U(0.5, 0.8)$;
jack placement jitter of $\pm$7\,cm on the table plane; jack splat
scale from $1.0$ to $1.3$; rendering gamma; grip point (45--55\,mm
behind the plug); and, at trajectory generation, grasp roll, jack
yaw, cable droop, lateral offsets, and approach distance. Key
calibration values are collected in Table~\ref{tab:nc-constants}.

\begin{table}[!htp]
\centering
\caption{Key simulation and rendering parameters. The connector
alignment and jack anchor encode corrected registration bugs (the
plug splat's insertion axis is local $+z$; the trajectory origin
sits 12\,mm past the jack mouth).}
\label{tab:nc-constants}
\resizebox{\columnwidth}{!}{%
\begin{tabular}{@{}lll@{}}
\toprule
\textbf{Quantity} & \textbf{Value} & \textbf{Note} \\
\midrule
Solver & VBD, 24 iterations, 8 substeps & 60\,Hz control \\
Contact & $k_e{=}10^{6}$, $k_d{=}0$, gap 0.1\,mm & SDF resolution 512 \\
Jack position & $(0.295, -0.876, 0.835)$\,m & world frame \\
Robot base & $(0.426, -0.106, 0.87)$\,m & \\
Jack anchor & $(0, 0, 0.018)$\,m & $0.030$ mouth $-$ $0.012$ seat-aim \\
Connector alignment & $(-90, 0, 0)^{\circ}$ RPY & insertion axis local $+z$ \\
Connector anchor & $(-1.5, -1.5, 17.2)$\,mm & mating face in plug splat \\
Grip angle & $-0.0152$\,rad & 6.35\,mm gap, 6.5\,mm cable \\
Camera & D415 $640\times480 \rightarrow 512^2$ & 30\,fps \\
\bottomrule
\end{tabular}%
}
\end{table}

 or paste before \end{document}.
% Needs only packages already in the preamble.

\providecolor{cmured}{RGB}{196, 18, 48}

%=====================================================================
\section{Simulation: Deformable Cable Insertion}
%=====================================================================

To study insertion with a deformable object, we simulate ethernet
cabling: the robot grips a flexible cable $\sim$5\,cm behind the
connector, so the RJ45 plug dangles from its strain-relief boot and
must be levelled and seated into a panel-mount jack with
sub-millimetre clearance. This section summarizes the design of the
simulation and the empirical lessons that shaped it.

\subsection{Choice of Physics Engine}

We evaluated MuJoCo, Genesis, Isaac Sim, and Newton. MuJoCo and
Genesis could not resolve the fine cable/connector interaction, which
left the two NVIDIA-stack options. We chose Newton over Isaac Sim
because its VBD solver treats one-dimensional deformables ---
self-collision, bend/twist transport, two-way coupling to rigid
bodies --- as a first-class, GPU-batched capability, and because it
is a headless Python library rather than a GUI-centric Omniverse
application, which kept iteration on contact behaviour fast. The cost
is maturity: PhysX articulations are far more battle-tested, and we
spent roughly fifteen debugging iterations discovering undocumented
VBD constraints (fixed joints freeze deformable chains unless
merged away; actuated joints must be softened; the rest pose must be
updated via forward kinematics before the solver is constructed;
stable joint-limit damping is $\sim 10^{-4}$, four orders of
magnitude below the conventional default).
Since Newton is also being integrated into Isaac Lab as a swappable
backend, we do not view the choice as locking us out of that
ecosystem.

\subsection{Cable and Contact Modelling}

We model the cable at two fidelities. For teacher training, the cable
is a Cosserat rod of 10\,mm capsule segments (radius 3.25\,mm,
spanning 6\,cm behind to 5\,cm ahead of the grasp), with the
connector attached as a mesh shape on the front segment ---
attaching it as a welded rigid body destabilizes VBD. A PPO teacher
trained in this environment reaches 94.4\,\% held-success at the
hardest curriculum stage. For large-scale data generation we
substitute a rigid approximation: a single free body carrying the
cord capsule, the plug mesh, and a small box at the grasp point that
gives the flat fingerpads a roll-lock a smooth capsule lacks.

Contact uses narrow-band signed-distance fields built directly on
the CAD meshes ($k_e = 10^{6}$, contact gap 0.1\,mm, socket
$\mu = 0.5$, cable $\mu = 2.0$), avoiding convex decomposition
entirely. Two findings dominated the design.
\textcolor{cmured}{First}, contact damping must be zero: the default
($k_d = 100$) acts through SDF shell contacts that fire at several
millimetres of actual separation, injecting up to
$\sim$400\,rad/s of spurious angular velocity into the light
(6.9\,g) connector. \textcolor{cmured}{Second}, contacts must never
act on articulation bodies: every grasp variant that allowed this
diverged. We therefore disable finger-mesh collision and realize the
grasp with two kinematic pad proxies placed at the measured pad
faces, closed to a $\sim$1.5\,mm effective squeeze on the 6.5\,mm
cable.

\subsection{Robot Model}

The arm (StandardBots RO1 with an AG-145 gripper) is driven
kinematically: end-effector deltas ($\pm$2\,mm, $\pm 0.75^{\circ}$
per 60\,Hz step) are converted to joint targets by
damped-least-squares IK and interpolated across the eight physics
substeps, so the shallow grasp is never violated within a step. For
hardware-faithful data we replace direct IK with the real controller
chain --- jerk-limited trajectory retiming, 500\,Hz streaming,
10\,ms transport delay, a 2\,ms plant --- which reaches 100\,\% held
success once
replanning is seeded with velocity and acceleration state (the naive
rest-to-rest replan advanced only 0.12\,\% of commanded motion).

\subsection{Digital-Twin Geometry}

Geometry comes from the manufacturer CAD models of the plug and jack
(McMaster-Carr 9953K216 and 1422N17). Our first conversion, a 0.10\,mm
voxel remesh, silently destroyed the task --- the real plug--cavity
clearance is $\sim$0.11\,mm, comparable to the voxel pitch, so the
gap quantized to zero. We instead tessellate the STEPs directly; the
narrow-band SDF tolerates the resulting non-watertight meshes and
preserves the true clearance. We also removed the latch catch ledge
from the physics model after finding that a rigid connector jams on
it (the real part inserts only because the plastic flexes). To keep
physics and rendered appearance consistent, a scan-matched variant
registers a photogrammetry scan of the real cable+plug into the CAD
frame (Kabsch alignment, sub-micron residual), so the collision mesh
matches the splat the cameras see. For real-world collection we
print a single-body fixture that holds the jack rigidly, with a
simulated twin registered to the socket to within 0.5\,mm.

\subsection{Determinism}

Kinematic quantities (joint streams, wrist poses) reproduce
bit-exactly across runs; contact dynamics do not, because VBD's
contact reduction uses GPU atomics and the resulting contact
ordering perturbs the friction impulse on the light connector ---
same-seed episodes near the success boundary can flip. We therefore
validate calibration transfer on kinematic channels (exact to
$2.5\times10^{-14}$ degrees) and evaluate dynamic outcomes
statistically.

%=====================================================================
\section{Photoreal Observations via Gaussian Splatting}
%=====================================================================

\subsection{Rendering Architecture and Scene}

Observations are produced by compositing the Newton state into a
Gaussian-splat reconstruction of the real scene. An external
renderer receives per-object world transforms and camera parameters
each frame and returns RGB; the interface is physics-agnostic, so
the same renderer serves our Isaac/PhysX pipeline. The scene comprises 19 splats: table, jack,
connector head and tail, and the 15 robot links.

Because a deformable cannot be represented by a rigid splat, the
cable itself is not rendered: the connector head and boot splats
ride the simulated plug-face pose, the cord behind them is simulated
but unrendered, and the plug splat is cropped to the rigid extent
the physics actually tracks. The plug splat is baked by
\emph{reusing the physics registration}, guaranteeing that pixels
and collision geometry cannot drift apart; the jack splat is
registered by point-pair correspondence, which also resolves the
arbitrary SfM scale of the reconstruction (0.8193; 0.91\,mm surface
residual). Robot splats keep only degree-0 spherical harmonics, as
the renderer misinterprets higher orders.

\subsection{Cameras}

All views use the intrinsics of the real Intel RealSense D415 at
$640\times480$ ($f_x{=}605.78$, $f_y{=}605.29$, $c_x{=}326.14$,
$c_y{=}229.88$), so simulated and real observations share a
projection model; both are square-resized to $512^2$ identically. Up
to four views are rendered --- front, side, mirror, and a dynamic
wrist camera whose pose and FOV are taken from the calibrated camera
mount in the robot model, so the eye-in-hand view is rigid by
construction. We
verify simulator--renderer alignment by rendering the Newton scene
from identical cameras and diffing against the GS frames; this
requires reproducing the wrist camera's $\sim$8$^{\circ}$ roll,
which the stock debug camera cannot represent.

\subsection{Shadows, Compositing, and Randomization}

Gaussians emit baked radiance, so dynamic objects cast no shadows
and appear pasted onto the scene. We recover shadows with a two-pass
path-traced ratio mask over the Newton geometry: a pass with and a
pass without dynamic occluders, whose luminance ratio (clamped at 1)
multiplies the GS frame. All terms except occlusion cancel between
the passes, so no segmentation is needed and the reconstruction's
own baked shadows are not double-darkened. Frames are then
gamma-corrected (1.8) and cleaned by Difix3D before dataset
conversion; at evaluation, live camera frames pass through the same
Difix so the policy sees its training distribution.

Domain randomization spans both halves of the pipeline: background
yaw over $U(0,360^{\circ})$ and shadow strength over $U(0.5, 0.8)$;
jack placement jitter of $\pm$7\,cm on the table plane; jack splat
scale from $1.0$ to $1.3$; rendering gamma; grip point (45--55\,mm
behind the plug); and, at trajectory generation, grasp roll, jack
yaw, cable droop, lateral offsets, and approach distance. Key
calibration values are collected in Table~\ref{tab:nc-constants}.

\begin{table}[!htp]
\centering
\caption{Key simulation and rendering parameters. The connector
alignment and jack anchor encode corrected registration bugs (the
plug splat's insertion axis is local $+z$; the trajectory origin
sits 12\,mm past the jack mouth).}
\label{tab:nc-constants}
\resizebox{\columnwidth}{!}{%
\begin{tabular}{@{}lll@{}}
\toprule
\textbf{Quantity} & \textbf{Value} & \textbf{Note} \\
\midrule
Solver & VBD, 24 iterations, 8 substeps & 60\,Hz control \\
Contact & $k_e{=}10^{6}$, $k_d{=}0$, gap 0.1\,mm & SDF resolution 512 \\
Jack position & $(0.295, -0.876, 0.835)$\,m & world frame \\
Robot base & $(0.426, -0.106, 0.87)$\,m & \\
Jack anchor & $(0, 0, 0.018)$\,m & $0.030$ mouth $-$ $0.012$ seat-aim \\
Connector alignment & $(-90, 0, 0)^{\circ}$ RPY & insertion axis local $+z$ \\
Connector anchor & $(-1.5, -1.5, 17.2)$\,mm & mating face in plug splat \\
Grip angle & $-0.0152$\,rad & 6.35\,mm gap, 6.5\,mm cable \\
Camera & D415 $640\times480 \rightarrow 512^2$ & 30\,fps \\
\bottomrule
\end{tabular}%
}
\end{table}

 or paste before \end{document}.
% Needs only packages already in the preamble.

\providecolor{cmured}{RGB}{196, 18, 48}

%=====================================================================
\section{Simulation: Deformable Cable Insertion}
%=====================================================================

To study insertion with a deformable object, we simulate ethernet
cabling: the robot grips a flexible cable $\sim$5\,cm behind the
connector, so the RJ45 plug dangles from its strain-relief boot and
must be levelled and seated into a panel-mount jack with
sub-millimetre clearance. This section summarizes the design of the
simulation and the empirical lessons that shaped it.

\subsection{Choice of Physics Engine}

We evaluated MuJoCo, Genesis, Isaac Sim, and Newton. MuJoCo and
Genesis could not resolve the fine cable/connector interaction, which
left the two NVIDIA-stack options. We chose Newton over Isaac Sim
because its VBD solver treats one-dimensional deformables ---
self-collision, bend/twist transport, two-way coupling to rigid
bodies --- as a first-class, GPU-batched capability, and because it
is a headless Python library rather than a GUI-centric Omniverse
application, which kept iteration on contact behaviour fast. The cost
is maturity: PhysX articulations are far more battle-tested, and we
spent roughly fifteen debugging iterations discovering undocumented
VBD constraints (fixed joints freeze deformable chains unless
merged away; actuated joints must be softened; the rest pose must be
updated via forward kinematics before the solver is constructed;
stable joint-limit damping is $\sim 10^{-4}$, four orders of
magnitude below the conventional default).
Since Newton is also being integrated into Isaac Lab as a swappable
backend, we do not view the choice as locking us out of that
ecosystem.

\subsection{Cable and Contact Modelling}

We model the cable at two fidelities. For teacher training, the cable
is a Cosserat rod of 10\,mm capsule segments (radius 3.25\,mm,
spanning 6\,cm behind to 5\,cm ahead of the grasp), with the
connector attached as a mesh shape on the front segment ---
attaching it as a welded rigid body destabilizes VBD. A PPO teacher
trained in this environment reaches 94.4\,\% held-success at the
hardest curriculum stage. For large-scale data generation we
substitute a rigid approximation: a single free body carrying the
cord capsule, the plug mesh, and a small box at the grasp point that
gives the flat fingerpads a roll-lock a smooth capsule lacks.

Contact uses narrow-band signed-distance fields built directly on
the CAD meshes ($k_e = 10^{6}$, contact gap 0.1\,mm, socket
$\mu = 0.5$, cable $\mu = 2.0$), avoiding convex decomposition
entirely. Two findings dominated the design.
\textcolor{cmured}{First}, contact damping must be zero: the default
($k_d = 100$) acts through SDF shell contacts that fire at several
millimetres of actual separation, injecting up to
$\sim$400\,rad/s of spurious angular velocity into the light
(6.9\,g) connector. \textcolor{cmured}{Second}, contacts must never
act on articulation bodies: every grasp variant that allowed this
diverged. We therefore disable finger-mesh collision and realize the
grasp with two kinematic pad proxies placed at the measured pad
faces, closed to a $\sim$1.5\,mm effective squeeze on the 6.5\,mm
cable.

\subsection{Robot Model}

The arm (StandardBots RO1 with an AG-145 gripper) is driven
kinematically: end-effector deltas ($\pm$2\,mm, $\pm 0.75^{\circ}$
per 60\,Hz step) are converted to joint targets by
damped-least-squares IK and interpolated across the eight physics
substeps, so the shallow grasp is never violated within a step. For
hardware-faithful data we replace direct IK with the real controller
chain --- jerk-limited trajectory retiming, 500\,Hz streaming,
10\,ms transport delay, a 2\,ms plant --- which reaches 100\,\% held
success once
replanning is seeded with velocity and acceleration state (the naive
rest-to-rest replan advanced only 0.12\,\% of commanded motion).

\subsection{Digital-Twin Geometry}

Geometry comes from the manufacturer CAD models of the plug and jack
(McMaster-Carr 9953K216 and 1422N17). Our first conversion, a 0.10\,mm
voxel remesh, silently destroyed the task --- the real plug--cavity
clearance is $\sim$0.11\,mm, comparable to the voxel pitch, so the
gap quantized to zero. We instead tessellate the STEPs directly; the
narrow-band SDF tolerates the resulting non-watertight meshes and
preserves the true clearance. We also removed the latch catch ledge
from the physics model after finding that a rigid connector jams on
it (the real part inserts only because the plastic flexes). To keep
physics and rendered appearance consistent, a scan-matched variant
registers a photogrammetry scan of the real cable+plug into the CAD
frame (Kabsch alignment, sub-micron residual), so the collision mesh
matches the splat the cameras see. For real-world collection we
print a single-body fixture that holds the jack rigidly, with a
simulated twin registered to the socket to within 0.5\,mm.

\subsection{Determinism}

Kinematic quantities (joint streams, wrist poses) reproduce
bit-exactly across runs; contact dynamics do not, because VBD's
contact reduction uses GPU atomics and the resulting contact
ordering perturbs the friction impulse on the light connector ---
same-seed episodes near the success boundary can flip. We therefore
validate calibration transfer on kinematic channels (exact to
$2.5\times10^{-14}$ degrees) and evaluate dynamic outcomes
statistically.

%=====================================================================
\section{Photoreal Observations via Gaussian Splatting}
%=====================================================================

\subsection{Rendering Architecture and Scene}

Observations are produced by compositing the Newton state into a
Gaussian-splat reconstruction of the real scene. An external
renderer receives per-object world transforms and camera parameters
each frame and returns RGB; the interface is physics-agnostic, so
the same renderer serves our Isaac/PhysX pipeline. The scene comprises 19 splats: table, jack,
connector head and tail, and the 15 robot links.

Because a deformable cannot be represented by a rigid splat, the
cable itself is not rendered: the connector head and boot splats
ride the simulated plug-face pose, the cord behind them is simulated
but unrendered, and the plug splat is cropped to the rigid extent
the physics actually tracks. The plug splat is baked by
\emph{reusing the physics registration}, guaranteeing that pixels
and collision geometry cannot drift apart; the jack splat is
registered by point-pair correspondence, which also resolves the
arbitrary SfM scale of the reconstruction (0.8193; 0.91\,mm surface
residual). Robot splats keep only degree-0 spherical harmonics, as
the renderer misinterprets higher orders.

\subsection{Cameras}

All views use the intrinsics of the real Intel RealSense D415 at
$640\times480$ ($f_x{=}605.78$, $f_y{=}605.29$, $c_x{=}326.14$,
$c_y{=}229.88$), so simulated and real observations share a
projection model; both are square-resized to $512^2$ identically. Up
to four views are rendered --- front, side, mirror, and a dynamic
wrist camera whose pose and FOV are taken from the calibrated camera
mount in the robot model, so the eye-in-hand view is rigid by
construction. We
verify simulator--renderer alignment by rendering the Newton scene
from identical cameras and diffing against the GS frames; this
requires reproducing the wrist camera's $\sim$8$^{\circ}$ roll,
which the stock debug camera cannot represent.

\subsection{Shadows, Compositing, and Randomization}

Gaussians emit baked radiance, so dynamic objects cast no shadows
and appear pasted onto the scene. We recover shadows with a two-pass
path-traced ratio mask over the Newton geometry: a pass with and a
pass without dynamic occluders, whose luminance ratio (clamped at 1)
multiplies the GS frame. All terms except occlusion cancel between
the passes, so no segmentation is needed and the reconstruction's
own baked shadows are not double-darkened. Frames are then
gamma-corrected (1.8) and cleaned by Difix3D before dataset
conversion; at evaluation, live camera frames pass through the same
Difix so the policy sees its training distribution.

Domain randomization spans both halves of the pipeline: background
yaw over $U(0,360^{\circ})$ and shadow strength over $U(0.5, 0.8)$;
jack placement jitter of $\pm$7\,cm on the table plane; jack splat
scale from $1.0$ to $1.3$; rendering gamma; grip point (45--55\,mm
behind the plug); and, at trajectory generation, grasp roll, jack
yaw, cable droop, lateral offsets, and approach distance. Key
calibration values are collected in Table~\ref{tab:nc-constants}.

\begin{table}[!htp]
\centering
\caption{Key simulation and rendering parameters. The connector
alignment and jack anchor encode corrected registration bugs (the
plug splat's insertion axis is local $+z$; the trajectory origin
sits 12\,mm past the jack mouth).}
\label{tab:nc-constants}
\resizebox{\columnwidth}{!}{%
\begin{tabular}{@{}lll@{}}
\toprule
\textbf{Quantity} & \textbf{Value} & \textbf{Note} \\
\midrule
Solver & VBD, 24 iterations, 8 substeps & 60\,Hz control \\
Contact & $k_e{=}10^{6}$, $k_d{=}0$, gap 0.1\,mm & SDF resolution 512 \\
Jack position & $(0.295, -0.876, 0.835)$\,m & world frame \\
Robot base & $(0.426, -0.106, 0.87)$\,m & \\
Jack anchor & $(0, 0, 0.018)$\,m & $0.030$ mouth $-$ $0.012$ seat-aim \\
Connector alignment & $(-90, 0, 0)^{\circ}$ RPY & insertion axis local $+z$ \\
Connector anchor & $(-1.5, -1.5, 17.2)$\,mm & mating face in plug splat \\
Grip angle & $-0.0152$\,rad & 6.35\,mm gap, 6.5\,mm cable \\
Camera & D415 $640\times480 \rightarrow 512^2$ & 30\,fps \\
\bottomrule
\end{tabular}%
}
\end{table}

 or paste before \end{document}.
% Needs only packages already in the preamble.

\providecolor{cmured}{RGB}{196, 18, 48}

%=====================================================================
\section{Simulation: Deformable Cable Insertion}
%=====================================================================

To study insertion with a deformable object, we simulate ethernet
cabling: the robot grips a flexible cable $\sim$5\,cm behind the
connector, so the RJ45 plug dangles from its strain-relief boot and
must be levelled and seated into a panel-mount jack with
sub-millimetre clearance. This section summarizes the design of the
simulation and the empirical lessons that shaped it.

\subsection{Choice of Physics Engine}

We evaluated MuJoCo, Genesis, Isaac Sim, and Newton. MuJoCo and
Genesis could not resolve the fine cable/connector interaction, which
left the two NVIDIA-stack options. We chose Newton over Isaac Sim
because its VBD solver treats one-dimensional deformables ---
self-collision, bend/twist transport, two-way coupling to rigid
bodies --- as a first-class, GPU-batched capability, and because it
is a headless Python library rather than a GUI-centric Omniverse
application, which kept iteration on contact behaviour fast. The cost
is maturity: PhysX articulations are far more battle-tested, and we
spent roughly fifteen debugging iterations discovering undocumented
VBD constraints (fixed joints freeze deformable chains unless
merged away; actuated joints must be softened; the rest pose must be
updated via forward kinematics before the solver is constructed;
stable joint-limit damping is $\sim 10^{-4}$, four orders of
magnitude below the conventional default).
Since Newton is also being integrated into Isaac Lab as a swappable
backend, we do not view the choice as locking us out of that
ecosystem.

\subsection{Cable and Contact Modelling}

We model the cable at two fidelities. For teacher training, the cable
is a Cosserat rod of 10\,mm capsule segments (radius 3.25\,mm,
spanning 6\,cm behind to 5\,cm ahead of the grasp), with the
connector attached as a mesh shape on the front segment ---
attaching it as a welded rigid body destabilizes VBD. A PPO teacher
trained in this environment reaches 94.4\,\% held-success at the
hardest curriculum stage. For large-scale data generation we
substitute a rigid approximation: a single free body carrying the
cord capsule, the plug mesh, and a small box at the grasp point that
gives the flat fingerpads a roll-lock a smooth capsule lacks.

Contact uses narrow-band signed-distance fields built directly on
the CAD meshes ($k_e = 10^{6}$, contact gap 0.1\,mm, socket
$\mu = 0.5$, cable $\mu = 2.0$), avoiding convex decomposition
entirely. Two findings dominated the design.
\textcolor{cmured}{First}, contact damping must be zero: the default
($k_d = 100$) acts through SDF shell contacts that fire at several
millimetres of actual separation, injecting up to
$\sim$400\,rad/s of spurious angular velocity into the light
(6.9\,g) connector. \textcolor{cmured}{Second}, contacts must never
act on articulation bodies: every grasp variant that allowed this
diverged. We therefore disable finger-mesh collision and realize the
grasp with two kinematic pad proxies placed at the measured pad
faces, closed to a $\sim$1.5\,mm effective squeeze on the 6.5\,mm
cable.

\subsection{Robot Model}

The arm (StandardBots RO1 with an AG-145 gripper) is driven
kinematically: end-effector deltas ($\pm$2\,mm, $\pm 0.75^{\circ}$
per 60\,Hz step) are converted to joint targets by
damped-least-squares IK and interpolated across the eight physics
substeps, so the shallow grasp is never violated within a step. For
hardware-faithful data we replace direct IK with the real controller
chain --- jerk-limited trajectory retiming, 500\,Hz streaming,
10\,ms transport delay, a 2\,ms plant --- which reaches 100\,\% held
success once
replanning is seeded with velocity and acceleration state (the naive
rest-to-rest replan advanced only 0.12\,\% of commanded motion).

\subsection{Digital-Twin Geometry}

Geometry comes from the manufacturer CAD models of the plug and jack
(McMaster-Carr 9953K216 and 1422N17). Our first conversion, a 0.10\,mm
voxel remesh, silently destroyed the task --- the real plug--cavity
clearance is $\sim$0.11\,mm, comparable to the voxel pitch, so the
gap quantized to zero. We instead tessellate the STEPs directly; the
narrow-band SDF tolerates the resulting non-watertight meshes and
preserves the true clearance. We also removed the latch catch ledge
from the physics model after finding that a rigid connector jams on
it (the real part inserts only because the plastic flexes). To keep
physics and rendered appearance consistent, a scan-matched variant
registers a photogrammetry scan of the real cable+plug into the CAD
frame (Kabsch alignment, sub-micron residual), so the collision mesh
matches the splat the cameras see. For real-world collection we
print a single-body fixture that holds the jack rigidly, with a
simulated twin registered to the socket to within 0.5\,mm.

\subsection{Determinism}

Kinematic quantities (joint streams, wrist poses) reproduce
bit-exactly across runs; contact dynamics do not, because VBD's
contact reduction uses GPU atomics and the resulting contact
ordering perturbs the friction impulse on the light connector ---
same-seed episodes near the success boundary can flip. We therefore
validate calibration transfer on kinematic channels (exact to
$2.5\times10^{-14}$ degrees) and evaluate dynamic outcomes
statistically.

%=====================================================================
\section{Photoreal Observations via Gaussian Splatting}
%=====================================================================

\subsection{Rendering Architecture and Scene}

Observations are produced by compositing the Newton state into a
Gaussian-splat reconstruction of the real scene. An external
renderer receives per-object world transforms and camera parameters
each frame and returns RGB; the interface is physics-agnostic, so
the same renderer serves our Isaac/PhysX pipeline. The scene comprises 19 splats: table, jack,
connector head and tail, and the 15 robot links.

Because a deformable cannot be represented by a rigid splat, the
cable itself is not rendered: the connector head and boot splats
ride the simulated plug-face pose, the cord behind them is simulated
but unrendered, and the plug splat is cropped to the rigid extent
the physics actually tracks. The plug splat is baked by
\emph{reusing the physics registration}, guaranteeing that pixels
and collision geometry cannot drift apart; the jack splat is
registered by point-pair correspondence, which also resolves the
arbitrary SfM scale of the reconstruction (0.8193; 0.91\,mm surface
residual). Robot splats keep only degree-0 spherical harmonics, as
the renderer misinterprets higher orders.

\subsection{Cameras}

All views use the intrinsics of the real Intel RealSense D415 at
$640\times480$ ($f_x{=}605.78$, $f_y{=}605.29$, $c_x{=}326.14$,
$c_y{=}229.88$), so simulated and real observations share a
projection model; both are square-resized to $512^2$ identically. Up
to four views are rendered --- front, side, mirror, and a dynamic
wrist camera whose pose and FOV are taken from the calibrated camera
mount in the robot model, so the eye-in-hand view is rigid by
construction. We
verify simulator--renderer alignment by rendering the Newton scene
from identical cameras and diffing against the GS frames; this
requires reproducing the wrist camera's $\sim$8$^{\circ}$ roll,
which the stock debug camera cannot represent.

\subsection{Shadows, Compositing, and Randomization}

Gaussians emit baked radiance, so dynamic objects cast no shadows
and appear pasted onto the scene. We recover shadows with a two-pass
path-traced ratio mask over the Newton geometry: a pass with and a
pass without dynamic occluders, whose luminance ratio (clamped at 1)
multiplies the GS frame. All terms except occlusion cancel between
the passes, so no segmentation is needed and the reconstruction's
own baked shadows are not double-darkened. Frames are then
gamma-corrected (1.8) and cleaned by Difix3D before dataset
conversion; at evaluation, live camera frames pass through the same
Difix so the policy sees its training distribution.

Domain randomization spans both halves of the pipeline: background
yaw over $U(0,360^{\circ})$ and shadow strength over $U(0.5, 0.8)$;
jack placement jitter of $\pm$7\,cm on the table plane; jack splat
scale from $1.0$ to $1.3$; rendering gamma; grip point (45--55\,mm
behind the plug); and, at trajectory generation, grasp roll, jack
yaw, cable droop, lateral offsets, and approach distance. Key
calibration values are collected in Table~\ref{tab:nc-constants}.

\begin{table}[!htp]
\centering
\caption{Key simulation and rendering parameters. The connector
alignment and jack anchor encode corrected registration bugs (the
plug splat's insertion axis is local $+z$; the trajectory origin
sits 12\,mm past the jack mouth).}
\label{tab:nc-constants}
\resizebox{\columnwidth}{!}{%
\begin{tabular}{@{}lll@{}}
\toprule
\textbf{Quantity} & \textbf{Value} & \textbf{Note} \\
\midrule
Solver & VBD, 24 iterations, 8 substeps & 60\,Hz control \\
Contact & $k_e{=}10^{6}$, $k_d{=}0$, gap 0.1\,mm & SDF resolution 512 \\
Jack position & $(0.295, -0.876, 0.835)$\,m & world frame \\
Robot base & $(0.426, -0.106, 0.87)$\,m & \\
Jack anchor & $(0, 0, 0.018)$\,m & $0.030$ mouth $-$ $0.012$ seat-aim \\
Connector alignment & $(-90, 0, 0)^{\circ}$ RPY & insertion axis local $+z$ \\
Connector anchor & $(-1.5, -1.5, 17.2)$\,mm & mating face in plug splat \\
Grip angle & $-0.0152$\,rad & 6.35\,mm gap, 6.5\,mm cable \\
Camera & D415 $640\times480 \rightarrow 512^2$ & 30\,fps \\
\bottomrule
\end{tabular}%
}
\end{table}

